\newpage
\chapter{Maintaining Order, Removal from Office, and Disciplinary Measures}

\section{Receiving a Consultation or Learning of a Problem}

The person responsible checks relevant posts or requests, their content and timing, and the accounts of the people involved.
Confirm the wishes of someone who seeks advice only. Unless harm is imminent, do not establish a violation on one person's account alone.
When referring to announced community rules, check their wording and the circumstances in which they apply.
Do not treat requests such as taking care with spoilers or choosing roles as grounds for discipline under Chapter 8, Article 4 of the Academy Regulations solely on that basis.
When considering disciplinary measures, identify specifically which ground in that Article applies.

\section{Ordinary Responses and Disciplinary Measures}

Once a problem has been confirmed, the person responsible considers whether contacting those involved or correcting facts will resolve it, and shares the matter with the other Deans as needed.
When imposing a disciplinary measure, check the person's account and record separately the facts falling under Chapter 8, Article 4 of the Academy Regulations and the type of measure under Article 5.
As a general rule, measures progress from warning to suspension of activities to expulsion. Apply Article 6 to a further violation in the same term after a warning, or during or after suspension.
Suspension of activities or expulsion without a prior warning is permitted only for an exceptionally serious violation under Article 7.
Inform the person of the established facts, grounds, disciplinary measure, its duration if applicable, and where to submit an appeal.
The person responsible records who participated in the decision, the notification, and the start and end times.

\section{Emergency Measures}

When harm to order or Student safety is occurring or imminent, temporary access restrictions or other measures may be taken to the extent necessary under Chapter 8, Article 8 of the Academy Regulations.
The person responsible records the scope and start time of the measure, notifies the person concerned of the reasons without delay, and shares the information with the other Deans.
The measure ceases to have effect no later than the last day of the academic term concerned.
Because this is not a disciplinary measure, assess separately under that chapter whether discipline is needed if a continuing restriction is necessary.

\section{Resignation, Removal from Office, and Appeals}

Officeholders other than the Director notify the Director of resignation; the Director notifies the Deans (Chapter 8, Article 2 of the Regulations).
For removal from office, check the grounds in Article 3. For an officeholder other than the Director, the Director makes the decision after consultation among the Deans; for the Director, the Deans decide through consultation.
Removal from office is not a disciplinary measure against the person as a Student. Consider discipline separately if necessary.
Inform the person of the removal and its reasons, and record the decision-maker and the handover of duties.

When receiving an appeal against removal from office, an emergency measure, or a disciplinary measure, Deans other than those who made the original decision review it under Chapter 8, Article 9 of the Regulations.
The reviewers check the facts, applicable provisions, and necessity of the measure, then inform the appellant of the conclusion and reasons and record them.
Handle requests to reconsider refusal of an application under Chapter 3 of this Handbook, separately from this appeal process.
